% from ega4-9-fr.tex
\oldpage[IV-3]{69}

\begin{proposition}[9.5.5]
\label{IV.9.5.5-fr}
Soient $f:X\to S$ un morphisme de présentation finie, $Z$ une partie
localement constructible de $X$ telle que, pour tout $s\in S$, $Z_s$ soit
localement fermé dans $X_s$. Pour tout $s\in S$, soit
$D(s)\subset\mathbf Z\cup\{-\infty\}$ l'ensemble des dimensions des
composantes irréductibles de $Z_s$. Alors la fonction $s\mapsto D(s)$ est
localement constructible dans $S$.
\end{proposition}

Soit $\Phi$ une partie finie de $\mathbf Z\cup\{-\infty\}$~; il s'agit de
montrer que l'ensemble des $s\in S$ tels que $D(s)=\Phi$ est localement
constructible~; compte tenu de (9.2.3), nous avons encore à vérifier les deux
conditions de (9.2.2, (iii)).

En ce qui concerne la première, il s'agit de voir que si $X$ est un préschéma
algébrique sur un corps $k$, $Z$ une partie localement fermée de $X$, $k'$ une
extension de $k$, $p:X_{(k')}\to X$ la projection canonique, alors l'ensemble
des dimensions des composantes irréductibles de $Z$ est le même que celui des
dimensions des composantes irréductibles de $p^{-1}(Z)$~; compte tenu de
l'existence d'un sous-préschéma de $X$ ayant $Z$ pour espace sous-jacent
(\textbf{I}, 5.2.1), cela résulte de (4.2.8).

Pour la seconde condition de (9.2.2, (iii)), on est dans le cas où $S$ est
noethérien et intègre de point générique $\eta$, et $f:X\to S$ est un
morphisme de type fini, de sorte que $X$ est noethérien. Le sous-espace
$Z_\eta$ de l'espace noethérien $X_\eta$ est par hypothèse localement fermé,
donc a un nombre fini de composantes irréductibles $Z_{i\eta}$, qui sont
localement fermées dans $X_\eta$. Il existe par suite pour chaque indice $i$
une partie localement fermée $Z_i$ de $X$ telle que $(Z_i)_\eta=Z_{i\eta}$,
donc si $Z'=\bigcup_i Z_i$, on a $Z'_\eta=Z_\eta$. Mais comme $Z$ et $Z'$ sont
localement constructibles, on peut, en remplaçant $S$ par un voisinage de
$\eta$, supposer que $Z'=Z$ (9.5.1). En outre, pour $i\neq j$,
$Z_{i\eta}\cap Z_{j\eta}$ est rare et \emph{fermé} dans $Z_{i\eta}$~; donc
(9.5.3 et 9.5.4), on peut encore supposer, en restreignant $S$, que pour
$j\neq i$, $(Z_i)_s\cap(Z_j)_s$ est rare et fermé dans $(Z_i)_s$. Comme $Z_i$
est localement fermé dans $X$, il y a un sous-préschéma de $X$ ayant $Z_i$
pour espace sous-jacent (et encore noté $Z_i$), qui est de type fini sur
$S$ (\textbf{I}, 6.3.5). Posons
$U_i=Z_i-\bigcup_{j\neq i}(Z_i\cap Z_j)$, qui est ouvert dans $Z_i$ et tel que,
pour tout $s\in S$, $(U_i)_s$ soit ouvert partout dense dans $(Z_i)_s$~; en
outre, par construction, les $U_i$ sont deux à deux sans point commun. Comme
$(Z_i)_s$ est un $k(s)$-préschéma algébrique, l'ensemble des dimensions des
composantes irréductibles de $Z_s=\bigcup_i(Z_i)_s$ est le même que l'ensemble
des dimensions des composantes irréductibles de la réunion des $(U_i)_s$
(4.1.1.3), chacune de ces composantes étant déjà une composante irréductible
de l'un des $(U_i)_s$. On peut donc se borner au cas où $Z=U_i$~; en outre,
comme $Z_\eta$ est alors irréductible, il n'y a qu'une seule composante
irréductible de $Z$ qui rencontre $X_\eta$ $(\mathbf 0_{\mathrm I},2.1.8)$ et
l'on peut évidemment, en restreignant $Z$, supposer $Z$ irréductible. On est
finalement ramené à prouver le cas particulier suivant de (9.5.5)~:

\begin{corollary}[9.5.6]
\label{IV.9.5.6-fr}
Soient $S$ un préschéma noethérien et irréductible de point générique $\eta$,
$X$ un préschéma irréductible, $f:X\to S$ un morphisme dominant de type fini.
Alors il existe un voisinage $U$ de $\eta$ dans $S$ tel que, pour tout
$s\in U$, toutes les composantes irréductibles de $X_s$ soient de dimension
$n=\dim(X_\eta)$.
\end{corollary}

On peut évidemment se borner au cas où $S=\operatorname{Spec}(A)$ est affine,
$A$ étant donc noethérien~; remplaçant $f$ par $f_{\mathrm{red}}$, qui est de
type fini (1.5.4, (vii)), on peut

